\documentclass[11pt]{article}

\usepackage[a4paper,top=0.9in,bottom=0.9in,left=0.95in,right=0.95in]{geometry}
\usepackage[T1]{fontenc}
\usepackage{lmodern}
\usepackage{microtype}
\usepackage{amsmath}
\usepackage{booktabs}
\usepackage{tabularx}
\usepackage{array}
\usepackage{enumitem}
\usepackage{xcolor}
\usepackage[hidelinks]{hyperref}
\usepackage{titlesec}

\titleformat{\section}{\large\bfseries}{\thesection}{0.6em}{}
\titleformat{\subsection}{\normalsize\bfseries}{\thesubsection}{0.5em}{}
\titlespacing*{\section}{0pt}{0.7\baselineskip}{0.25\baselineskip}
\titlespacing*{\subsection}{0pt}{0.5\baselineskip}{0.15\baselineskip}

\setlist{topsep=2pt,itemsep=1pt,parsep=0pt}

\newcolumntype{Y}{>{\raggedright\arraybackslash}X}
\setlength{\emergencystretch}{3em}
\hbadness=10000

\title{\vspace{-1.2cm}Bitcoin Wallet Risk Analyzer\\[2pt]\large Test Plan and Evaluation Document}
\author{Ori Sinvani \textnormal{(325770824)} \quad Harel Brener \textnormal{(214179012)} \quad Nehoray Chalfon \textnormal{(325833531)}\\[3pt]
\normalsize\normalfont Project Advisor: Ofer Zevin}
\date{May 2026}

\begin{document}
\maketitle
\vspace{-0.6cm}

\section{Project Summary}

The goal of this project is to decide, given a Bitcoin wallet address, whether it has been used for criminal
activity or for benign activity. The classifier is a graph neural network that reads a small one-hop neighbourhood
around the wallet: the wallet itself sits at the centre with twelve behavioural features describing its lifetime,
how often it transacts, its fee patterns, and the size and direction of the amounts it sends and receives. The
counter-parties the wallet has interacted with form the surrounding nodes, and each interaction becomes a directed
edge carrying its log amount, its direction, and a normalised timestamp. A three-layer GATv2 graph attention
network with multi-head attention, residual connections, and a hybrid readout aggregates this neighbourhood into a
single risk score. Two real-world labelled datasets, REAL-CATS and Elliptic++, are merged into a balanced corpus
that is split once into a training set and a held-out test set.

The headline result on the held-out test split is an $F_1$ of about 92\% for the GNN, against about 84\% for an
XGBoost classifier trained on the same twelve tabular features without the graph context. The roughly eight-point
gap is the central finding: a one-hop view of the surrounding transaction structure carries information that the
tabular features alone do not capture.

\section{Experimental Data}\label{sec:data}

\paragraph{Sources.}
\begin{itemize}
  \item \textbf{REAL-CATS} (Su et al., 2024). An expert-curated dataset of Bitcoin wallets linked to criminal
        activity, compiled by domain analysts from public criminal-wallet reports (sanctions lists, ransomware
        disclosures, law-enforcement notices, and similar sources), paired with a benign control set. The labels
        come from human reporting rather than from on-chain behavioural inference. The dataset was released in
        2024, with its underlying material collected through roughly 2023 to 2024.
  \item \textbf{Elliptic++} (Elmougy and Liu, KDD 2023). An extension of the original Elliptic dataset (Weber et
        al., 2019). The underlying Bitcoin transactions span roughly 2014 to 2017 and are organised into 49
        time-steps of about two weeks each. The dataset was released in 2023.
\end{itemize}

\paragraph{Unified feature matrix.}
Both datasets are mapped onto a common schema of nineteen features that can be computed from either source: eleven
shared base features (total transactions, total sent and received amounts, total fees, lifetime, send and receive
counts, max and min sent and received), three additional features that can be calculated for both datasets (mean
blocks between transactions, mean fee share, average transaction size), and five engineered features
(send/receive ratio, activity rate, fee per transaction, transaction-size range, and in/out balance). Every
feature is also stored as its $\log(1+x)$ transform. After correlation pruning, in which seven features that were
highly correlated with another feature ($r > 0.95$) were dropped, the model uses the twelve least-redundant
log-scaled features.

\paragraph{Balanced split.}
A balancing step drops cross-source addresses whose labels disagree, collapses duplicates whose labels agree, takes
every criminal wallet and a matched random sample of benign wallets within each source, and then produces a
stratified 80/20 train/test split keyed on the joint label and source label. The result is that both classes and
both data sources are represented in identical proportions in the training set and in the test set. The random
seed is fixed throughout.

\begin{table}[h]
\centering\small
\begin{tabular}{lrrrrr}
\toprule
Split & Rows & Criminal & Benign & REAL-CATS & Elliptic++ \\
\midrule
Training set & 86{,}870 & 43{,}435 & 43{,}435 & 64{,}044 & 22{,}826 \\
Test set     & 21{,}718 & 10{,}859 & 10{,}859 & 16{,}012 &  5{,}706 \\
\bottomrule
\end{tabular}
\end{table}

\paragraph{Per-wallet ego-graphs.}
For each wallet a small graph is built. The wallet sits at the centre carrying the twelve log features. Every
counter-party it has transacted with appears as a ghost node, and at forward time the model replaces the
zero-feature placeholder of each ghost node with a learned embedding plus a node-type embedding so that ghost
neighbours are not silently treated as fresh wallets with all-zero behaviour. Each interaction is encoded as a
directed edge with three features: the log of the amount transferred, a direction flag (1 for outgoing, 0 for
incoming), and a timestamp normalised against the Bitcoin genesis block.

\section{Baseline Models}

\begin{itemize}
  \item \textbf{XGBoost (tabular).} Trained on the same training set with the same twelve features and a standard
        scaler, evaluated on the same test set. The contrast with the GNN therefore isolates the value of the
        one-hop neighbourhood structure: same features, same split, only the graph view differs.
  \item \textbf{Basic GCN (graph baseline).} A vanilla graph convolutional network applied to the same ego-graphs,
        without our model's architectural additions: no learnable ghost-node embedding, no input-to-final residual
        connection, and a plain mean-pool readout instead of the hybrid readout. The contrast with the GNN
        isolates the value of those additions, given that the underlying graph view is the same.
\end{itemize}

\section{Evaluation Metrics}

The model is judged primarily on $F_1$ on the held-out test split, complemented by the standard set of
classification metrics so that a single number does not hide the full picture.

\begin{table}[h]
\centering\small
\begin{tabularx}{\linewidth}{l Y}
\toprule
\textbf{Family} & \textbf{Metric} \\
\midrule
Discrimination     & Accuracy, Precision, Recall, $F_1$, ROC-AUC, PR-AUC \\
Per-class          & Precision, Recall, and $F_1$ for the benign class and the criminal class separately \\
Confusion          & Confusion matrix on the test set \\
Calibration        & Brier score, Expected Calibration Error (ECE), and a reliability diagram \\
\bottomrule
\end{tabularx}
\end{table}

\noindent Model selection during training uses the validation $F_1$, with early stopping. Beyond the table above,
the evaluation step also dumps the misclassified test addresses to CSV (one file for false positives, one for
false negatives) and writes the per-feature gradient-saliency scores to JSON. These files are saved for future
analysis.

\section{Experimental Setting}

\paragraph{Train, validation, and test split.}
The 80/20 stratified split described in Section~\ref{sec:data} produces the canonical training and test sets.
Inside training, the training set is further partitioned 85/15 into a training pool and a validation pool using a
fixed seed, so the validation pool is identical across runs and is used both for early stopping and for picking
the best epoch.

\paragraph{Calibration split.}
A small slice is carved out of the training pool before the train/validation split and is reserved exclusively for
post-hoc temperature scaling. An LBFGS optimiser fits a single temperature parameter on the negative log-likelihood
of this slice, and the resulting scalar is saved alongside the model and applied at inference time in both the
backend and the analysis notebook.

\paragraph{Train/inference parity.}
At inference time, when a live wallet is analysed, the same feature pipeline is run on the raw transaction data
that was used to build the training matrix: satoshi-to-BTC scaling, $\log(1+x)$ transform, and the same clipping
ranges that were applied during dataset construction (activity rate clipped to $[0, 1000]$, send/receive ratio and
in/out balance clipped to $[0, 100]$, fee share clipped to $[0, 1]$, and transaction-size range computed from
sent-only amounts).

\paragraph{Reproducibility.}
A single random seed (42) is used for the dataset split, for the validation sub-split, for the dataloader shuffle,
and for the torch RNG. Every evaluation output can be regenerated from the saved model weights and calibration
scalar.

\section{Parameters}\label{sec:params}

\begin{table}[h]
\centering\small
\begin{tabularx}{\linewidth}{l Y}
\toprule
\textbf{Component} & \textbf{Setting} \\
\midrule
GNN architecture     & 3 $\times$ GATv2Conv layers (4, 4, and 2 attention heads), hidden size 64, edge dim 3 \\
                     & dropout 0.2, final dropout 0.3, DropEdge rate 0.1 \\
                     & hybrid readout: centre node $\oplus$ mean-pool $\oplus$ max-pool ($3 \times 64 = 192$) \\
                     & classifier $192 \rightarrow 64 \rightarrow 2$ \\
Optimiser            & AdamW, learning rate $10^{-3}$, weight decay $0.01$, gradient clip max-norm $1.0$ \\
Scheduler            & OneCycleLR, max learning rate $5 \times 10^{-3}$, pct-start $0.1$, cosine anneal \\
Loss                 & Cross-entropy, label smoothing $0$ \\
Schedule             & up to 150 epochs, batch size 64, early stop on validation $F_1$, patience 20 \\
Calibration          & LBFGS, max-iter 50, learning rate $0.01$ \\
XGBoost              & n-estimators $= 200$, max-depth $= 8$, lr $= 0.1$, subsample $= 0.8$, colsample-bytree $= 0.8$ \\
Basic GCN            & 3 $\times$ GCNConv layers, hidden size 64, mean-pool readout, dropout 0.2 \\
\bottomrule
\end{tabularx}
\end{table}

\section{Test Plan: Detailed Breakdown}\label{sec:tests}

The tests are grouped into two blocks: \textit{model quality}, which checks how well the model performs and
generalises against the two baselines, and \textit{model behaviour}, which probes what it has actually learned.
Each test produces a JSON or CSV file that can be reviewed after the run.

\subsection*{Model quality}

\begin{description}[leftmargin=0pt,labelindent=0pt,style=nextline,itemsep=4pt,parsep=0pt]
  \item[\textit{Held-out test evaluation.}] Run the evaluation pipeline end to end and report Accuracy, Precision,
    Recall, $F_1$, ROC-AUC, and PR-AUC for our GNN, for the XGBoost tabular baseline, and for the Basic GCN
    graph baseline. The two contrasts isolate the value of the graph view (XGBoost vs.\ GNN) and the value of
    our architectural additions on top of a plain GCN (Basic GCN vs.\ GNN).
  \item[\textit{Per-source generalisation.}] Restrict the same metrics to the REAL-CATS-only and Elliptic++-only
    subsets of the test split, to make sure neither source is carrying the headline number on its own.
  \item[\textit{Calibration quality.}] Brier score, Expected Calibration Error, and a reliability diagram on the
    held-out test set after temperature scaling, so that the predicted probabilities can be used as risk
    thresholds downstream.
  \item[\textit{Live-wallet evaluation.}] Build a small evaluation set of roughly 200 live wallet addresses
    (sampled outside the training distribution) and measure precision at the top-$k$ predicted-criminal scores.
    This is the strongest available test of generalisation to fresh, real-world traffic.
\end{description}

\subsection*{Model behaviour}

\begin{description}[leftmargin=0pt,labelindent=0pt,style=nextline,itemsep=4pt,parsep=0pt]
  \item[\textit{Misclassification dump.}] Export the false positives and false negatives on the held-out test
    set to two CSV files, sorted by predicted probability. Saved for future analysis.
  \item[\textit{Interpretability.}] Gradient-saliency feature importance over the twelve input features; rank
    the top drivers for each class and check whether they are stable across the criminal and benign cohorts.
\end{description}

\section{Schedule and Timeline}

The work runs over seven weeks, ending with the final submission on \textbf{25 June 2026}.

\begin{table}[h]
\centering\small
\begin{tabularx}{\linewidth}{l l Y}
\toprule
\textbf{Phase} & \textbf{Theme} & \textbf{Work} \\
\midrule
Week 1 & Data completion       & Finish the remaining graph downloads (from $\sim$38k to the full 86k training graphs); meeting with the project advisor. \\
Week 2 & Model retraining      & Retrain the GNN on the full training set. \\
Week 3 & Live-wallet evaluation & Build the live evaluation set ($\sim$200 addresses) and measure precision at top-$k$. \\
Week 4 & Deployment            & Deploy the backend and frontend to the university server. \\
Week 5 & Release and report    & Write the detailed evaluation report. \\
Week 6 & Final prep            & Build the final slide deck. \\
Week 7 & Buffer and polish     & Polish, dry runs, and final review before submission. \\
\bottomrule
\end{tabularx}
\end{table}

\end{document}
